% TODO make hw stylesheet
\documentclass[11pt]{scrartcl}
\usepackage[a4paper, total={6in, 8in}]{geometry}
\usepackage[usenames,dvipsnames,svgnames]{xcolor}
\usepackage[shortlabels]{enumitem}
\usepackage[framemethod=TikZ]{mdframed}
\usepackage{amsmath,amssymb,amsthm}
\usepackage[colorlinks]{hyperref}
\usepackage{multicol}
\usepackage{tabularx}
\usepackage{setspace}
\setstretch{1.05}

\hypersetup{
	colorlinks=true,
	linkcolor=blue,
	filecolor=magenta,      
	urlcolor=magenta,
	pdftitle={MAT 319 HW}
}

\usepackage{microtype}
\usepackage{mathtools}
\usepackage[headsepline]{scrlayer-scrpage}
\usepackage{thmtools}
\usepackage{todonotes}

\addtolength{\textheight}{3.14cm}
\providecommand{\ol}{\overline}
\providecommand{\eps}{\varepsilon}
\providecommand{\half}{\frac{1}{2}}
\providecommand{\dang}{\measuredangle} %% Directed angle
\providecommand{\CC}{\mathbb C}
\providecommand{\FF}{\mathbb F}
\providecommand{\NN}{\mathbb N}
\providecommand{\QQ}{\mathbb Q}
\providecommand{\RR}{\mathbb R}
\providecommand{\ZZ}{\mathbb Z}
\providecommand{\dg}{^\circ}
\providecommand{\ii}{\item}
\providecommand{\alert}[1]{{\sffamily\textbf{\textcolor{blue}{#1}}}}

\reversemarginpar

\theoremstyle{remark}
\newtheorem{innercustomthm}{Problem}
\newenvironment{problem}[1]
{\renewcommand\theinnercustomthm{\textit{#1}}\innercustomthm}
{\endinnercustomthm}

\newenvironment{soln}{\begin{proof}[Solution]}{\end{proof}}
\newenvironment{walkthrough}{\noindent \textbf{Walkthrough.}}{\medskip}
\newlist{walk}{enumerate}{3}
\setlist[walk]{label=\bfseries (\alph*)}

\providecommand{\pow}{\operatorname{Pow}}
\providecommand{\vocab}[1]{{\textbf{\textcolor{ForestGreen}{#1}}}}
\providecommand{\lcm}{\operatorname{lcm}}
\providecommand{\abs}[1]{\left|#1\right|}

\usepackage{asymptote}
\begin{asydef}
	defaultpen(fontsize(10pt));
	unitsize(1cm);
	size(8cm); // set a reasonable default
	usepackage("amsmath");
	usepackage("amssymb");
	settings.tex="pdflatex";
	settings.outformat="pdf";
	// Replacement for olympiad+cse5 which is not standard
	import geometry;
	// recalibrate fill and filldraw for conics
	void filldraw(picture pic = currentpicture, conic g, pen fillpen=defaultpen, pen drawpen=defaultpen)
	{ filldraw(pic, (path) g, fillpen, drawpen); }
	void fill(picture pic = currentpicture, conic g, pen p=defaultpen)
	{ filldraw(pic, (path) g, p); }
	// some geometry
	pair foot(pair P, pair A, pair B) { return foot(triangle(A,B,P).VC); }
	pair orthocenter(pair A, pair B, pair C) { return orthocentercenter(A,B,C); }
	pair centroid(pair A, pair B, pair C) { return (A+B+C)/3; }
	// cse5 abbreviations
	path CP(pair P, pair A) { return circle(P, abs(A-P)); }
	path CR(pair P, real r) { return circle(P, r); }
	pair IP(path p, path q) { return intersectionpoints(p,q)[0]; }
	pair OP(path p, path q) { return intersectionpoints(p,q)[1]; }
	path Line(pair A, pair B, real a=0.6, real b=a) { return (a*(A-B)+A)--(b*(B-A)+B); }
	// cse5 more useful functions
	picture CC() {
		picture p=rotate(0)*currentpicture;
		currentpicture.erase();
		return p;
	}
	pair MP(Label s, pair A, pair B = plain.S, pen p = defaultpen) {
		Label L = s;
		L.s = "$"+s.s+"$";
		label(L, A, B, p);
		return A;
	}
	pair Drawing(Label s = "", pair A, pair B = plain.S, pen p = defaultpen) {
		dot(MP(s, A, B, p), p);
		return A;
	}
	path Drawing(path g, pen p = defaultpen, arrowbar ar = None) {
		draw(g, p, ar);
		return g;
	}
\end{asydef}

\usepackage{mathrsfs}

\ihead{\footnotesize MAT 319 HW02}
\ohead{\footnotesize \today}

\title{\vspace{-2em}MAT 319 HW02}
\author{\large Aditya Pahuja}
\date{\large Due 9/11}
\begin{document}
	\maketitle
	
	\section{Textbook Problems}
	\begin{problem}{7.1}
		Write out the first five terms of the following sequences.
		\begin{enumerate}[\textbf{(\alph*)}]
			\ii $s_n = \frac{1}{3n + 1}$
			\ii $b_n = \frac{3n + 1}{4n - 1}$
			\ii $c_n = \frac{n}{3^n}$
			\ii $\sin(\frac{n\pi}{4})$
		\end{enumerate}
	\end{problem}
	\begin{soln}
		Part (a):
		$\frac{1}{4},\;\frac{1}{7},\;\frac{1}{10},\;\frac{1}{13},\;\frac{1}{16}$
		
		Part (b):
		$\frac{4}{3},\;\frac{7}{7},\;\frac{10}{11},\;\frac{13}{15},
		\;\frac{16}{19}$
		
		Part (c):
		$\frac{1}{3},\;\frac{2}{9},\;\frac{3}{27},\;\frac{4}{81},
		\;\frac{5}{243}$
		
		Part (d):
		$\frac{\sqrt2}{2},\;1,\;\frac{\sqrt2}{2},\;0,\;-\frac{\sqrt2}{2}$
	\end{soln}
	
	\begin{problem}{7.3}
		For each sequence below, determine whether it converges and,
		if it converges, give its limit. No proofs are required.
		
		\begin{tabularx}{\linewidth}{XX}
			\textbf{(a)} $a_n = \frac{n}{n + 1}$& \textbf{(b)} $b_n = \frac{n^2 + 3}{n^2 - 3}$\\
			\textbf{(c)} $c_n = 2^{-n}$ & \textbf{(d)} $t_n = 1 + \frac 2n$\\
			\textbf{(e)} $x_n = 73 + (-1)^n$ & \textbf{(f)} $s_n = 2^{1/n}$\\
			\textbf{(g)} $y_n = n!$& \textbf{(h)} $d_n = (-1)^nn$\\
			\textbf{(i)} $\frac{(-1)^n}{n}$& \textbf{(j)} $\frac{7n^3 + 8n}{2n^3 - 3}$\\
			\textbf{(k)} $\frac{9n^2 - 18}{6n + 18}$& \textbf{(l)} $\sin(\frac{n\pi}{2})$\\
			\textbf{(m)} $\sin(n\pi)$& \textbf{(n)} $\sin(\frac{2n\pi}{3})$\\
			\textbf{(o)} $\frac{\sin n}{n}$& \textbf{(p)} $\frac{2^{n+1} + 5}{2^n - 7}$\\
			\textbf{(q)} $\frac{3^n}{n!}$& \textbf{(r)} $(1 + \frac 1n)^2$\\
			\textbf{(s)} $\frac{4n^2 + 3}{3n^2 - 2}$& \textbf{(t)} $\frac{6n + 4}{9n^2 + 7}$
		\end{tabularx}
	\end{problem}
	
	\begin{soln}\phantom{buh}
		
		\begin{tabularx}{\linewidth}{XX}
			\textbf{(a)} Converges to $1$. & \textbf{(b)} Converges to $1$.\\
			\textbf{(c)} Converges to $0$. & \textbf{(d)} Converges to $1$.\\
			\textbf{(e)} Does not converge.& \textbf{(f)} Converges to $1$.\\
			\textbf{(g)} Does not converge (diverges to $+\infty$).& \textbf{(h)} Does not converge.\\
			\textbf{(i)} Converges to $0$.& \textbf{(j)} Converges to $\frac 72$.\\
			\textbf{(k)} Does not converge (diverges to $+\infty$).& \textbf{(l)} Does not converge.\\
			\textbf{(m)} Converges to $0$.& \textbf{(n)} Does not converge.\\
			\textbf{(o)} Converges to $0$.& \textbf{(p)} Converges to $2$.\\
			\textbf{(q)} Converges to $0$.& \textbf{(r)} Converges to $1$.\\
			\textbf{(s)} Converges to $\frac43$.& \textbf{(t)} Converges to $0$.
		\end{tabularx}
	\end{soln}
	
	\begin{problem}{8.1}
		Prove the following:
		\begin{enumerate}[\textbf{(\alph*)}]
			\ii $\lim \frac{(-1)^n}{n} = 0$
			\ii $\lim \frac{1}{n^{1/3}} = 0$
			\ii $\lim \frac{2n - 1}{3n + 2} = \frac23$
		\end{enumerate}
	\end{problem}
	
	\begin{soln}
		For part $(a)$, we want to show that, for every $\eps > 0$,
		\[\abs{\frac{(-1)^n}{n} - 0} < \eps\]
		is eventually true. We show that as long as $n > \frac1\eps$,
		the inequality holds; indeed,
		\[\abs{\frac{(-1)^n}{n} - 0} = \abs{\frac{(-1)^n}{n}} = \frac 1n
		< \frac{1}{1/\eps} = \eps,\]
		so the inequality is eventually true and the limit of the sequence is $0$.
		
		For part $(b)$, we want to show that, for every $\eps > 0$,
		\[\abs{\frac{1}{n^{1/3}} - 0} < \eps\]
		is eventually true. We show that as long as $n > \frac1{\eps^3}$,
		the inequality holds; indeed,
		\[\abs{\frac{1}{n^{1/3}} - 0} = \abs{\frac{1}{n^{1/3}}} = \frac 1{n^{1/3}}
		< \frac{1}{(1/\eps^3)^{1/3}} = \frac{1}{1/\eps} = \eps,\]
		so the inequality is eventually true and the limit of the sequence is $0$.
		
		For part $(c)$, we want to show that, for every $\eps > 0$,
		\[\abs{\frac{2n-1}{3n+2} - \frac23} < \eps\]
		% (2n - 1)/(3n + 2) - 2/3 = [(6n - 3) - (6n + 4)]/(9n + 6) = 7/(9n + 6) < eps
		% 7/eps - 6 < n
		is eventually true. We show that as long as $n > \frac{7}{9\eps} - \frac23$,
		the inequality holds; indeed,
		\begin{align*}
			\abs{\frac{2n-1}{3n+2} - \frac23} &=
			\abs{\frac{3(2n - 1) - 2(3n + 2)}{3(3n + 2)}}\\ &=
			\abs{\frac{-7}{9n + 6}} = \frac{7}{9n + 6}\\ &<
			\frac{7}{9\left(\frac{7}{9\eps} - \frac23\right) + 6}\\
			&= \frac{7}{7/\eps - 6 + 6} = \eps,
		\end{align*}
		so the inequality is eventually true and the limit of the sequence is
		$\frac23$.
	\end{soln}
	
	\begin{problem}{8.2}
		Determine, with proof, the limits of the following sequences.
		\begin{enumerate}[\textbf{(\alph*)}]
			\ii $a_n = \frac{n}{n^2 + 1}$
			\ii $b_n = \frac{7n - 19}{3n + 7}$
			\ii $c_n = \frac{4n + 3}{7n - 5}$
		\end{enumerate}
	\end{problem}
	\begin{soln}
		The limit of $a_n$ is $\boxed 0$.
		To prove this, we want to show that, for all $\eps > 0$,
		\[|a_n - 0| < \eps\]
		is eventually true. I claim the inequality holds whenever $n > \frac1\eps$.
		Indeed, the condition implies
		\[\abs{\frac{n}{n^2 + 1} - 0} = \frac{n}{n^2 + 1} < \frac{n}{n^2}
		= \frac{1}{n} < \frac{1}{1/\eps} = \eps,\]
		so the inequality is eventually true and thus the limit is $0$.
		
		The limit of $b_n$ is \framebox{$\frac73$}.
		To prove this, we want to show that, for all $\eps > 0$,
		\[\abs{b_n - \frac73} < \eps\]
		is eventually true. I claim the inequality holds whenever
		$n > \frac{106}{9\eps} - \frac{7}{3}$.
		Indeed, the condition implies
		\begin{align*}
			\abs{\frac{7n - 19}{3n + 7} - \frac73} &=
			\abs{\frac{3(7n - 19) - 7(3n + 7)}{3(3n + 7)}}\\
			&= \abs{\frac{-106}{9n + 21}} = \frac{106}{9n + 21}\\
			&< \frac{106}{9\left(\frac{106}{9\eps} - \frac{7}{3}\right) + 21}\\
			&= \frac{106}{106/\eps - 21 + 21} = \eps,
		\end{align*}
		so the inequality is eventually true and thus the limit is $\frac73$.
		
		The limit of $c_n$ is \framebox{$\frac47$}.
		To prove this, we want to show that, for all $\eps > 0$,
		\[\abs{c_n - \frac47} < \eps\]
		is eventually true. I claim the inequality holds whenever
		$n > \frac{43}{49\eps} + \frac{5}{7}$.
		Indeed, the condition implies
		\begin{align*}
			\abs{\frac{4n + 3}{7n - 5} - \frac47} &=
			\abs{\frac{7(4n + 3) - 4(7n - 5)}{7(7n-5)}}\\
			&= \abs{\frac{43}{49n - 35}} = \frac{43}{49n - 35}\\
			&< \frac{43}{49\left(\frac{43}{49\eps} + \frac{5}{7}\right) - 35}\\
			&= \frac{43}{43/\eps + 35 - 35} = \eps,
		\end{align*}
		so the inequality is eventually true and thus the limit is $\frac47$.
	\end{soln}
	
	\begin{problem}{9.13}
		Show that
		\[\lim_{n\to\infty} a^n =\begin{cases}
			0 & \text{if }|a| < 1\\
			1 & \text{if }a = 1\\
			+\infty & \text{if }a > 1\\
			\text{does not exist} & \text{if }a\le -1.
		\end{cases}\]
	\end{problem}
	\begin{soln}
		Clearly if $a = 1$ then $a^n = 1^n = 1$ for all $n$,
		so that case is obvious.
		
		If $a > 1$, then we want to show that $\lim_{n\to\infty} a^n = +\infty$,
		which means that, for each $M\in\RR$, $a^n$ is eventually always greater
		than $M$. To show this, write $a = 1 + x$ for a positive real $x$. Then,
		\[a^n =
		(1 + x)^n = 1 + nx + \binom n2 x^2 + \binom n3 x^3 + \cdots + x^n > 1 + nx\]
		by the binomial theorem. This means that, if $n > \frac{M - 1}{x}$,
		\[a^n > 1 + nx = M,\]
		so $a^n$ is eventually greater than $M$, meaning $(a, a^2, a^3, \dots)$
		diverges to $+\infty$.
		
		Now, we treat the case of $|a| < 1$ by splitting it into the three cases
		$a = 0$, $0 < a < 1$, and $-1 < a < 0$.
		\begin{itemize}
			\ii If $a = 0$, then $a^n = 0$ for all $n$, so the limit is $0$
			(because $\eps > 0$ means that $\eps > 0 = \abs{a^n - 0}$).
			\ii If $0 < a < 1$, suppose for the sake of contradiction
			that $\lim_{n\to\infty} a^n$ is nonzero, say equal to some
			positive real $c$ ($c$ must be positive because $a^n > 0$ for all $n$).
			Then, for all $\eps > 0$, the inequality $|a^n - c| < \eps$
			is eventually true, say for all $n > N$. This means that for $n > N$,
			$a^n < c + \eps$ because
			\[c + \eps > |c| + |a^n - c| \ge |a^n| = a^n,\]
			so, for all such $n$,
			\[\abs{\frac{1}{a^n} - \frac{1}{c}} =
			\abs{\frac{c - a^n}{ca^n}} < \frac{\eps}{c(c + \eps)}
			< \frac{\eps}{c^2}\]
			where the first inequality comes from $|c - a^n| < \eps$
			and $\frac{1}{c + \eps} < \frac{1}{a^n}$.
			I claim that this means
			\[\lim_{n\to\infty} \frac{1}{a^n} = \frac{1}{c}.\]
			To show this, we want to prove that for all $\eps' > 0$,
			\[\abs{\frac{1}{a^n} - c} < \eps'.\]
			We showed above that, for any $\eps > 0$,
			\[\abs{\frac{1}{a^n} - \frac1c} > \frac{\eps}{c^2},\]
			so picking $\eps = c^2\eps'$ shows that there is an $N'$
			such that $n > N'$ implies
			\[\abs{\frac{1}{a^n} - c} < \eps',\]
			which is enough to show that the limit of $\frac{1}{a^n}$
			as $n$ goes to infinity is $\frac1c$.
			This, finally, is our contradiction:
			we have shown that the limit of $(\frac1a)^n$ is $\frac1c$,
			which means the sequence given by $(\frac1a)^n$ should converge;
			however, we showed earlier that, since $\frac1a > 1$,
			this sequence diverges to $+\infty$! Thus, $\lim_{n\to\infty} a^n = 0$,
			as desired.
			\ii For the last case, $-1 < a < 0$, we want to show that,
			for each $\eps > 0$,
			\[\abs{a^n - 0} = \abs{a^n} = |a|^n < \eps\]
			is eventually true.
			This is equivalent to showing that
			\[\abs{|a|^n - 0} < \eps\]
			is eventually true --- that is, it suffices to show that
			the limit as $n$ goes to infinity of $|a|^n$ is zero.
			This follows directly from the case where $0 < a < 1$
			because $|a|$ is between $0$ and $1$, so we are done!
		\end{itemize}
		
		Finally, we need to show that $a < -1$ implies that
		$\lim_{n\to\infty} a^n$ doesn't exist. Suppose for the sake of contradiction
		that the limit does exist. Then, the limit is either a real number
		or $\pm\infty$.
		In the former case, let the limit be $L$.
		This means that for all $\eps > 0$,
		\[|a^n - L| < \eps\]
		is eventually true, so it must be true for $\eps = \frac12$.
		We find by the triangle inequality that
		\[1 = 2\eps > |a^n - L| + |a^{n+1} - L| = |a^n - L| + |L - a^{n+1}| \ge
		|a^n - a^{n+1}| = |a|^n + |a|^{n+1}\]
		with the last equality following from $a < 0$, since $a^n$
		and $a^{n+1}$ have different signs.
		Then, since $|a| > 1$, $|a|^m > 1$ for all $m$, which means
		\[1 = 2\eps > |a|^n + |a|^{n+1} > 1 + 1 = 2,\]
		which is a contradiction.
		
		Therefore, if the limit exists,
		it must be one of $+\infty$ and $-\infty$. This means either
		for all real $M$, $a^n > M$ eventually, or for all real $M$,
		$a^n < M$ eventually. Taking $M = 0$ shows that statement is false,
		since $a^n < M$ if $n$ is odd and $a^n > M$ if $n$ is even,
		so $a^n < M$ and $a^n > M$ are both frequently true
		and thus both are also frequently false
		(meaning neither is eventually true).
		This shows that the limit doesn't exist if $a < -1$.
	\end{soln}
	
	\pagebreak
	\begin{problem}{9.18}\phantom{buh}
		\begin{enumerate}[\textbf{(\alph*)}]
			\ii Verify $1 + a + a^2 + \cdots + a^n = \frac{a^{n+1} - 1}{a - 1}$
			for all $a\ne 1$.
			\ii Find $\lim_{n\to\infty} (1 + a + a^2 + \cdots + a^n)$ for $|a| < 1$.
			\ii Calculate
			$\lim_{n\to\infty}(1 + \frac13 + \frac19 + \cdots + \frac{1}{3^n})$.
			\ii What is $\lim_{n\to\infty} (1 + a + a^2 + \cdots + a^n)$
			for $a\ge1$?
		\end{enumerate}
	\end{problem}
	\begin{soln}
		For part (a), we can see that
		\begin{align*}
			(a - 1)(1 + a + a^2 + \cdots + a^n) &=
			a(1 + a + a^2 + \cdots + a^n) - (1 + a + a^2 + \cdots + a^n)\\
			&= (a + a^2 + a^3 + \cdots + a^{n+1}) - (1 + a + a^2 + \cdots + a^n)
		\end{align*}
		and then $a$, $a^2$, \dots, $a^n$ all cancel out,
		leaving behind just $a^{n+1} - 1$; dividing by $a - 1$
		shows the desired equality.
		(Alternatively, induction also works nicely.)
		
		For part (b), I claim that
		\[\lim_{n\to\infty} (1 + a + a^2 + \cdots + a^n) = \boxed{\frac{1}{1 - a}}\]
		for $|a| < 1$. To see this, we simplify the expression in the limit
		via part (a), at which point it suffices by the definition of limit
		to prove that, for all $\eps > 0$,
		\[\abs{\frac{a^{n + 1} - 1}{a - 1} - \frac{1}{1 - a}} < \eps\]
		is eventually true. This is fairly easy to prove:
		it ends up being equivalent to show that
		\[\eps > \abs{\frac{a^{n + 1} - 1}{a - 1} - \frac{1}{1 - a}}
		= \abs{\frac{a^{n+1}}{a - 1}}
		\iff |a - 1|\eps > |a^{n+1}| = a^{n+1}\]
		is eventually true. We know from Problem 9.13 that the limit of $a^n$
		as $n$ goes to infinity is $0$ when $|a| < 1$, so, for each $\eps' > 0$,
		there exists an $N$ such that $n > N$ implies $|a^n - 0| = a^n > \eps'$.
		Taking $\eps' = |a - 1|\eps$ shows that
		\[|a - 1|\eps > a^{n+1}\]
		is true for $n > N$ (noting that $n + 1 > n > N$),
		i.e. the desired inequality is eventually true.
		Therefore, the limit is indeed $\frac{1}{1 - a}$.
		
		Part (c) is just part (b) in the case where $a = \frac13$
		(note that $|\frac13| < 1$, which is the necessary condition for
		the expression in part (b) to work), so the desired limit is just
		\[\frac{1}{1 - \frac13} = \boxed{\frac32}.\]
		
		Finally, for part (d), the answer is $+\infty$. We see that $a \ge 1$
		implies $a^n\ge 1$ for all $a$ --- this is easy to show by e.g. induction,
		since $a^n \ge 1$ implies $a^{n+1} = a^n\cdot a \ge 1$. Therefore,
		\[1 + a + a^2 + \cdots + a^n \ge n\]
		for all $n$. This means that, for each real $M$, if $N$ is the smallest
		natural number greater than $M$, then $n > N$ implies
		\[1 + a + a^2 + \cdots + a^n \ge n > N > M,\]
		which is exactly what it means for the limit to be $+\infty$.
	\end{soln}
	
	\section{The Other Problems}
	\begin{problem}{1}
		Say if the following statements are true or false.
		Carefully justify your answer.
		\begin{enumerate}[(1)]
			\ii A statement $P(n)$, $n\in\NN$, is true frequently if and only if
			it is true for infinitely many values of $n$.
			\ii $\sqrt n\in\NN$ frequently.
			\ii $\sqrt n\in\QQ$ frequently.
		\end{enumerate}
	\end{problem}
	\begin{soln}
		(1) is true. We prove the equivalent statement that
		$P(n)$ is not true frequently (i.e. is eventually false)
		if and only if $P(n)$ is true for finitely many $n$.
		
		First, if $P(n)$ is eventually false, then by definition there is an $N$
		such that $n > N$ implies $P(n)$ is false. This means that if $P(n)$ is true
		then $n \le N$, but, since there are only finitely many natural numbers
		less than $N$, $P(n)$ can only be true for finitely many $n$.
		
		Conversely, if $P(n)$ is true for finitely many $n$,
		then the set of $n$ for which $P(n)$ is true has some largest value $N$.
		In particular, $n > N$ implies $P(n)$ is false, i.e. $P(n)$
		is eventually false as desired.
		
		(2) and (3) are both true. We just prove (2), since
		$\sqrt n\in \NN$ frequently implies $\sqrt n\in\QQ$ frequently
		because $\NN\subseteq \QQ$. To prove (2), part (1) says that
		we just need to find an infinite set of numbers
		whose square roots are positive integers;
		the obvious set $\{1^2, 2^2, 3^2, 4^2, \dots\}$ does the job.
	\end{soln}
	
	\begin{problem}{2}
		For $a_n = n(n + 1)$, say if the following statements are true or false.
		\begin{enumerate}[(1)]
			\ii $a_n$ is eventually even,
			\ii $a_n$ is eventually odd,
			\ii $a_n$ is frequently even,
			\ii $a_n$ is frequently odd.
		\end{enumerate}
	\end{problem}
	\begin{soln}
		\framebox{(1) and (3) are true, (2) and (4) are false}.
		
		We begin by observing that $n$ and $n + 1$ are consecutive integers,
		so one of them must be odd and the other must be even
		regardless of the value of $n$. This means that $n(n + 1)$
		is a product of an odd and an even integer, so it is always even.
		
		Since $a_n$ is even for all $n \ge 1$, it is trivially
		eventually even, which also means it is frequently even, so (1) and (3)
		are true.
		
		Additionally, (2) and (4) are false because $a_n$ being
		frequently odd would require there to be infinitely many odd terms
		in $(a_n)$ (which there are obviously not), and $a_n$
		not being frequently odd implies that it is not eventually odd, either.
	\end{soln}
	\pagebreak
	\begin{problem}{3}
		Let $a_n$ be some sequence.
		Say if the following statements are true or false, and justify your answer.
		\begin{enumerate}[(1)]
			\ii If $a_n\to +\infty$ then $a_n$ must be eventually increasing.
			\ii If $a_n > 0$ and $a_n$ is unbounded then $a_n$ must diverge to
			$+\infty$.
			\ii If $a_n\to 0$ then either $a_n\to 0^+$ or $a_n\to 0^-$.
			\ii If $a_n\to 0^+$ then $a_n$ must be eventually decreasing.
		\end{enumerate}
	\end{problem}
	
	\begin{soln}
		(1) is false: consider the sequence $(a_n)$ given by
		$(2, 1, 4, 3, 6, 5, 8, 7, \dots)$, which alternates between
		the smallest positive even number not currently in the sequence
		and the smallest positive odd number not currently in the sequence.
		This is clearly unbounded, since for any $M$, $n > 100!\cdot|M| + 10^{10}$
		implies $a_n > M$ (in particular $a_n > n - 1$ for all $n$,
		so the bound on $n$ is slightly overkill).
		
		(2) is false: take $b_n = 2^n$ if $n$ is even and $1$ if $n$ is odd.
		Then, all terms of the sequence are positive and, since
		$\lim_{k\to\infty} 2^k = +\infty$, $b_n$ must be unbounded (otherwise
		the powers of two would have an upper bound).
		However, this only means that, for each real $M$,
		$b_n$ is \emph{frequently} greater than $M$;
		for the sequence to diverge to $+\infty$, it would need to be
		\emph{eventually} greater than $M$, which does not hold if, say, $M = \pi$
		(any $M\ge 1$ works as a counterexample).
		
		(3) is false: take $c_n = \frac{(-1)^n}{n}$.
		This alternates between being positive and negative,
		but its limit is $0$, as shown in textbook problem 8.1(a).
		
		(4) is false: take $d_n = \frac 1{a_n}$
		(where $(a_n)$ is the sequence $(2, 1, 4, 3, 6, 5, \dots)$
		defined earlier). For each $\eps$, $|d_n - 0| = c_n = \frac{1}{a_n} < \eps$
		is eventually true because it holds whenever $n > \frac{100!}{\eps} + 1000$,
		since $a_n > n - 1$ implies
		\[d_n < \frac{1}{n - 1} = \frac{1}{\frac{100!}{\eps} + 999}
		< \frac{\eps}{100!} < \eps.\]
		This means that $c_n\to 0^+$, but the sequence is not eventually
		decreasing because
		$\frac{1}{2k} < \frac{1}{2k-1}$ for all positive integers $k$.
	\end{soln}
	
\end{document}